\begin{center}
\textbf{\LARGE Autorisation de diffusion et d’archivage du rapport}
\end{center}

\subsection*{Informations de l’auteur}
\begin{itemize}[label={--}]
    \item NOM et Prénom de l’auteur : \underline{\hspace{5cm}}
    \item Spécialité : \underline{\hspace{5cm}}
    \item Titre du projet : \underline{\hspace{5cm}}
    \item Nom de l’entreprise : \underline{\hspace{5cm}}
    \item NOM et Prénom du tuteur du Projet de Fin d'Études : \underline{\hspace{5cm}}
    \item NOM et Prénom du correspondant pédagogique INSA : \underline{\hspace{5cm}}
\end{itemize}

\section*{Archivage du rapport de PFE}
À l'issue de son stage, l'étudiant(e) stagiaire rédigera un rapport qui devra être communiqué aussi bien à l'organisme d'accueil qu'à l'établissement d'enseignement supérieur pour évaluation.  
Par obligation réglementaire (instruction 2005-003 parue au B.O du 16 juin 2005), l’INSA de Rennes est tenu de conserver pour archive une version du rapport de PFE. Tout rapport de PFE sera donc systématiquement archivé par le département de l'auteur de l’INSA de Rennes.   
Si le rapport contient des données confidentielles, une version expurgée pourra être archivée à la place de la version intégrale.

\begin{itemize}
    \item ☐ archivage de la version intégrale
    \item ☐ archivage de la version expurgée, merci de justifier : \underline{\hspace{5cm}}
\end{itemize}

\section*{Diffusion du rapport de PFE}
\subsection*{Autorisation de diffusion par l’entreprise commanditaire}
Nous, soussignés, représentant de l’entreprise commanditaire :  
\begin{itemize}[label={--}]
    \item Nom : \underline{\hspace{5cm}}
    \item Prénom : \underline{\hspace{5cm}}
    \item Fonction : \underline{\hspace{5cm}}
\end{itemize}

\begin{itemize}
    \item ☐ autorisons le signalement et la diffusion du document désigné ci-dessus sur l’intranet du département de l'auteur de l’INSA
    \item ☐ autorisons uniquement le signalement du document désigné ci-dessus sur l’intranet du département de l'auteur de l’INSA
    \item ☐ n’autorisons ni le signalement ni la diffusion du document désigné ci-dessus sur l’intranet du département de l'auteur de l’INSA et demandons la confidentialité (5 ans maximum), jusqu’à la date suivante (mois/année) : \underline{\hspace{3cm}}
\end{itemize}

Fait à \underline{\hspace{3cm}}, le \underline{\hspace{3cm}} \\
Signature et/ou cachet de l’entreprise commanditaire : \underline{\hspace{5cm}}

\subsection*{Autorisation de diffusion par l’auteur}
Je soussigné(e), \\
Nom : \underline{\hspace{5cm}} \\
Prénom : \underline{\hspace{5cm}}

\begin{itemize}
    \item ☐ autorise le signalement et la diffusion du document désigné ci-dessus sur l’intranet du département de l'auteur de l’INSA
    \item ☐ autorise uniquement le signalement du document désigné ci-dessus sur l’intranet du département de l'auteur de l’INSA
    \item ☐ n’autorise ni le signalement ni la diffusion du document désigné ci-dessus sur l’intranet du département de l'auteur de l’INSA
\end{itemize}

Fait à \underline{\hspace{3cm}}, le \underline{\hspace{3cm}} \\
Signature : \underline{\hspace{5cm}}

\subsection*{Autorisation de diffusion par le correspondant pédagogique de l’INSA}
Par défaut, le correspondant pédagogique, membre de l’équipe enseignante, est supposé avoir une contribution mineure sur le rapport et ne pas s’opposer à sa diffusion. Merci de contacter le responsable des stages du département en cas d’opposition.


